\documentclass[11pt,a4paper]{article}

\usepackage[utf8]{inputenc}
\usepackage[T1]{fontenc}
\usepackage{lmodern}
\usepackage[french]{babel}
\usepackage[margin=2.1cm]{geometry}
\usepackage{microtype}
\usepackage{amsmath,amssymb,amsthm}
\usepackage{booktabs}
\usepackage{array}
\usepackage[table]{xcolor}
\usepackage{graphicx}
\usepackage{caption}
\usepackage[most]{tcolorbox}

\definecolor{navy}{HTML}{1F3864}
\definecolor{bluemid}{HTML}{2E5496}
\definecolor{lightblue}{HTML}{D6E0F0}

\captionsetup{font=small,labelfont={bf,color=navy},labelsep=period}
\graphicspath{{../figures/}}

\newtcolorbox{cle}[1][]{colback=lightblue!40,colframe=navy,boxrule=0.6pt,
  arc=2pt,left=8pt,right=8pt,top=6pt,bottom=6pt,#1}

\theoremstyle{definition}
\newtheorem{proposition}{Proposition}
\newtheorem{definition}[proposition]{Définition}

\newcommand{\T}{^{\top}}
\newcommand{\E}{\mathbb{E}}

\title{\vspace{-1.2cm}\color{navy}\bfseries Le socle mathématique de $W=S+A$\\[2pt]
  \large Flux stationnaire, renversement du temps, fluctuation martingale
  et non-identifiabilité de la direction}
\author{Kélian Kaddouri}
\date{30 juillet 2026}

\begin{document}
\maketitle
\vspace{-0.9cm}

\begin{cle}
\textbf{Ce que cette fiche démontre.} La fiche \emph{réversibilité et martingalisation}
énonce trois identités et les vérifie numériquement. Celle-ci les \textbf{démontre}, et en
tire le résultat qui porte le chapitre~9~: une statistique construite sur des
co-occurrences non ordonnées prend la même valeur sur un processus et sur son renversé,
donc elle ne peut pas identifier la direction de la contagion. Le placebo $z=-0{,}33$
cesse d'être un échec de calibration pour devenir une \textbf{conséquence de la structure
de l'observable}. Et pourtant la direction déplace le capital~: c'est exactement la
configuration qui appelle une identification partielle.
\end{cle}

% =====================================================================
\section*{1. Le bon objet n'est pas $P$, c'est le flux stationnaire}

Une remarque préalable sans laquelle l'argument est faux. On ne décompose pas la matrice de
transition $P$ elle-même~: une matrice stochastique symétrique force une loi stationnaire
uniforme, si bien que la symétrie de $P$ n'a pas de contenu probabiliste. L'objet correct est
le \textbf{flux stationnaire}
\begin{equation}
  F_{jk} \;=\; \pi_j P_{jk},
  \label{eq:flux}
\end{equation}
c'est-à-dire la probabilité, en régime stationnaire, d'observer la transition $j \to k$
sur un pas de temps. C'est lui qui se coupe en deux, et c'est lui que les données agrégées
approchent.

On note $\|M\|^2=\sum_{jk}M_{jk}^2$ la norme de Frobenius et
$\langle M,N\rangle=\sum_{jk}M_{jk}N_{jk}$ le produit scalaire associé.

% =====================================================================
\section*{2. La décomposition, et son unicité}

\begin{proposition}[Décomposition orthogonale]\label{prop:ortho}
Toute matrice réelle $F$ s'écrit de manière unique $F=S+A$ avec $S\T=S$ et $A\T=-A$,
à savoir $S=\tfrac12(F+F\T)$ et $A=\tfrac12(F-F\T)$. Ces deux parties sont orthogonales,
$\langle S,A\rangle=0$, et par conséquent
\begin{equation}
  \|F\|^2 \;=\; \|S\|^2 + \|A\|^2 .
  \label{eq:pythagore}
\end{equation}
\end{proposition}

\begin{proof}
L'existence est le calcul direct $S+A=F$. Pour l'unicité, si $F=S'+A'$ avec $S'$ symétrique
et $A'$ antisymétrique, alors $F\T=S'-A'$, d'où $S'=\tfrac12(F+F\T)$ et
$A'=\tfrac12(F-F\T)$. Pour l'orthogonalité, on renomme les indices puis on utilise les deux
symétries~:
\[
  \langle S,A\rangle=\sum_{jk}S_{jk}A_{jk}=\sum_{jk}S_{kj}A_{kj}
  =\sum_{jk}S_{jk}(-A_{jk})=-\langle S,A\rangle,
\]
donc $\langle S,A\rangle=0$. L'identité~\eqref{eq:pythagore} suit en développant
$\|S+A\|^2$.
\end{proof}

En composantes, avec le \textbf{courant de probabilité}
$J_{jk}=\pi_jP_{jk}-\pi_kP_{kj}$~:
\begin{equation}
  S_{jk}=\tfrac12\bigl(\pi_jP_{jk}+\pi_kP_{kj}\bigr),
  \qquad
  A_{jk}=\tfrac12 J_{jk}.
  \label{eq:composantes}
\end{equation}
$S_{jk}$ est le \textbf{trafic} entre les piliers $j$ et $k$, $A_{jk}$ est le \textbf{sens}
dans lequel il penche. L'orthogonalité dit que ces deux informations ne se recouvrent pas~:
connaître l'intensité n'apprend rien sur la direction, et réciproquement. C'est déjà
l'annonce de ce qui suit.

% =====================================================================
\section*{3. Le renversement du temps retourne $A$ et laisse $S$}

\begin{proposition}[Chaîne adjointe]\label{prop:renversement}
Soit $(X_t)$ stationnaire de loi $\pi>0$. Le processus renversé $(X_{T-t})$ est une chaîne
de Markov de matrice de transition
\begin{equation}
  P^{*}_{jk}\;=\;\frac{\pi_k P_{kj}}{\pi_j},
  \label{eq:adjointe}
\end{equation}
de même loi stationnaire $\pi$, et son flux vérifie
\begin{equation}
  F^{*} \;=\; F\T \;=\; S - A .
  \label{eq:reversal}
\end{equation}
En particulier la chaîne est réversible si et seulement si $A=0$, condition qui s'écrit
$\pi_jP_{jk}=\pi_kP_{kj}$ pour tout couple $(j,k)$ \emph{(bilan détaillé)}.
\end{proposition}

\begin{proof}
$P^{*}$ est bien stochastique~: $\sum_k P^{*}_{jk}=\pi_j^{-1}\sum_k\pi_kP_{kj}=\pi_j^{-1}
\pi_j=1$ par stationnarité de $\pi$. Son flux vaut $\pi_jP^{*}_{jk}=\pi_kP_{kj}=F_{kj}$,
d'où $F^{*}=F\T$, et $F\T=(S+A)\T=S-A$. Enfin $P^{*}=P$ équivaut à $F=F\T$, donc à $A=0$.
\end{proof}

\begin{cle}
\textbf{L'identité centrale, en une phrase.} Renverser le temps laisse $S$ \emph{intacte}
et retourne le \emph{signe} de $A$. $S$ est la part du processus qui ignore dans quel sens
le temps s'écoule~; $A$ est la seule part qui le sait.
\end{cle}

% =====================================================================
\section*{4. La fluctuation martingale ne voit que $S$}

C'est le point que le parallèle avec la martingalisation devait établir. Soit $L=P-I$ le
générateur et $f$ une fonction sur les cinq piliers. La \textbf{martingale de Dynkin}
\begin{equation}
  M^f_t \;=\; f(X_t)-f(X_0)-\sum_{s<t}(Lf)(X_s)
  \label{eq:dynkin}
\end{equation}
est la fluctuation pure~: on a retranché exactement la dérive prévisible. C'est cela, la
martingalisation. Sa variation quadratique prévisible s'écrit
$\langle M^f\rangle_t=\sum_{s<t}\Gamma(f,f)(X_s)$ avec le carré du champ
$\Gamma(f,f)(j)=\sum_k P_{jk}\bigl(f(k)-f(j)\bigr)^2$.

\begin{proposition}[La fluctuation est aveugle à la direction]\label{prop:dirichlet}
En régime stationnaire, l'intensité moyenne de la fluctuation ne dépend que de $S$~:
\begin{equation}
  \E_\pi\bigl[\Gamma(f,f)\bigr]
  \;=\;\sum_{j,k}S_{jk}\bigl(f(k)-f(j)\bigr)^2
  \;=\;2\,\mathcal{E}(f,f),
  \label{eq:dirichlet}
\end{equation}
où $\mathcal{E}$ est la forme de Dirichlet. La part antisymétrique $A$ y contribue
exactement zéro, pour toute fonction test $f$.
\end{proposition}

\begin{proof}
On développe et on décompose~:
\[
  \E_\pi[\Gamma(f,f)]=\sum_j\pi_j\sum_kP_{jk}(f(k)-f(j))^2
  =\sum_{j,k}F_{jk}(f(k)-f(j))^2
  =\sum_{j,k}(S_{jk}+A_{jk})(f(k)-f(j))^2 .
\]
Le poids $(f(k)-f(j))^2$ est \emph{symétrique} en $(j,k)$ tandis que $A$ est
\emph{antisymétrique}~: en appariant le terme $(j,k)$ avec le terme $(k,j)$, la somme
$\sum_{j,k}A_{jk}(f(k)-f(j))^2$ se compense terme à terme et vaut $0$. Il reste
l'expression annoncée.
\end{proof}

\medskip
\noindent\textbf{La version en temps continu est encore plus nette.} Pour une diffusion
$dX_t=b(X_t)\,dt+\sigma\,dB_t$, la partie martingale de $f(X_t)$ vaut
$\int_0^t\nabla f\cdot\sigma\,dB_s$, de crochet
\begin{equation}
  \bigl\langle M^f\bigr\rangle_t=\int_0^t\bigl|\sigma\T\nabla f(X_s)\bigr|^2\,ds .
  \label{eq:crochet-diffusion}
\end{equation}
Le crochet ne contient \textbf{que} le coefficient de diffusion $\sigma$~: la dérive $b$
n'y apparaît pas du tout. La fluctuation voit la diffusion, jamais la partie rotationnelle
de la dérive. C'est le même énoncé que la proposition~\ref{prop:dirichlet}, dans l'autre
cadre.

% =====================================================================
\section*{5. Le résultat de non-identifiabilité}

\begin{definition}[Statistique invariante par renversement]\label{def:invariante}
Une statistique $T$ est \emph{invariante par renversement} si $T(F)=T(F\T)$ pour toute
matrice de flux $F$.
\end{definition}

\begin{proposition}[Non-identifiabilité de la direction]\label{prop:nonid}
Soit $T$ invariante par renversement. Alors $T(S+A)=T(S-A)$ pour tout $A$ antisymétrique.
Par conséquent la loi de $T$ sous le paramètre $(S,A)$ est identique à sa loi sous
$(S,-A)$~: la vraisemblance est \textbf{plate} dans la direction de $A$, et $A$ n'est pas
identifié par la famille des lois de $T$.
\end{proposition}

\begin{proof}
Par la proposition~\ref{prop:renversement}, $(S+A)\T=S-A$. L'invariance donne alors
$T(S+A)=T\bigl((S+A)\T\bigr)=T(S-A)$. Deux valeurs distinctes du paramètre engendrant la
même loi des observations, l'identifiabilité au sens usuel est en défaut.
\end{proof}

\begin{cle}[colback=orange!8,colframe=orange!70!black]
\textbf{L'application au placebo, et le changement de statut du chapitre~9.} Une base
annuelle agrégée ne retient que des co-occurrences \textbf{non ordonnées}~: on y observe
que deux piliers ont flambé la même année, jamais lequel a précédé l'autre. Toute
statistique construite sur cette base est par construction invariante par renversement.
La proposition~\ref{prop:nonid} interdit donc qu'elle détecte $A$.

Le placebo $z=-0{,}33$ n'est ni un manque de puissance ni un problème de taille
d'échantillon, et \textbf{aucun volume de données annuelles ne le réparera}. On passe de
\og je n'ai pas réussi à calibrer la direction\fg{} à \og je démontre que la direction
n'est pas calibrable sur cette classe de données\fg{}, ce qui n'est pas la même phrase
devant un jury.
\end{cle}

% =====================================================================
\section*{6. Un scalaire pour mesurer l'irréversibilité}

\begin{proposition}[Production d'entropie]\label{prop:entropie}
La quantité
\begin{equation}
  \sigma\;=\;\frac12\sum_{j,k}\bigl(F_{jk}-F_{kj}\bigr)\,
  \ln\!\frac{F_{jk}}{F_{kj}}
  \label{eq:entropie}
\end{equation}
vérifie $\sigma\ge 0$, avec égalité si et seulement si $A=0$.
\end{proposition}

\begin{proof}
Chaque terme est de la forme $(a-b)\ln(a/b)$ avec $a,b>0$. Si $a>b$ les deux facteurs sont
positifs, si $a<b$ ils sont tous deux négatifs, et le produit est nul exactement lorsque
$a=b$. La somme de termes positifs est nulle si et seulement si chacun l'est, c'est-à-dire
$F_{jk}=F_{kj}$ pour tout couple, soit $A=0$.
\end{proof}

L'interprétation est la plus parlante des trois. $\sigma$ est le taux d'entropie relative
entre la loi de la trajectoire lue à l'endroit et celle de la trajectoire lue à l'envers~:
\begin{equation}
  \sigma\;=\;\lim_{T\to\infty}\frac1T\,
  \mathrm{KL}\Bigl(\mathcal{L}(X_0,\dots,X_T)\ \Big\|\ \mathcal{L}(X_T,\dots,X_0)\Bigr).
  \label{eq:kl}
\end{equation}
Autrement dit $\sigma$ mesure \textbf{à quel point le film à l'endroit est distinguable du
film à l'envers}. Si $\sigma=0$, aucun observateur, si bien équipé soit-il, ne peut dire
dans quel sens le temps s'écoule, donc aucun ne peut reconstituer la direction. La limite
d'identification n'est pas une faiblesse de méthode, c'est une contrainte de principe.

\medskip
\noindent\textbf{Le critère de Kolmogorov}, version vérifiable de la même chose~: la chaîne
est réversible si et seulement si, pour tout cycle $j_1\to j_2\to\cdots\to j_n\to j_1$,
\begin{equation}
  P_{j_1j_2}P_{j_2j_3}\cdots P_{j_nj_1}
  \;=\;
  P_{j_1j_n}\cdots P_{j_3j_2}P_{j_2j_1},
  \label{eq:kolmogorov}
\end{equation}
c'est-à-dire si le produit le long de la boucle est le même dans les deux sens. Lorsque les
deux produits diffèrent, il existe une \textbf{circulation} véritable. La direction de la
contagion est donc une circulation dans le graphe des cinq piliers, et une donnée agrégée,
qui ne conserve que des marges et des co-occurrences, ne voit jamais les cycles.

% =====================================================================
\section*{7. Pourquoi $A$ compte malgré tout pour le capital}

Il faut énoncer ce point, faute de quoi on objectera que si $A$ est invisible, il est sans
enjeu. C'est l'inverse.

\begin{proposition}[$A$ n'est pas décoratif]\label{prop:trace}
Pour $W=S+A$ on a $\langle S,A\rangle=0$ et
\begin{equation}
  \operatorname{tr}\!\bigl(W^2\bigr)\;=\;\|S\|^2-\|A\|^2 .
  \label{eq:trace}
\end{equation}
\end{proposition}

\begin{proof}
$\operatorname{tr}(W^2)=\operatorname{tr}(S^2)+2\operatorname{tr}(SA)+
\operatorname{tr}(A^2)$. Le terme croisé est nul, car
$\operatorname{tr}(SA)=\langle S\T,A\rangle=\langle S,A\rangle=0$ par la
proposition~\ref{prop:ortho}. Ensuite $\operatorname{tr}(S^2)=\|S\|^2$ puisque $S$ est
symétrique, et $\operatorname{tr}(A^2)=-\operatorname{tr}(AA\T)=-\|A\|^2$ puisque
$A\T=-A$.
\end{proof}

Le capital passe par l'inverse de Leontief,
$(I-W)^{-1}=I+W+W^2+W^3+\cdots$, et dès le terme quadratique
$W^2=S^2+SA+AS+A^2$ fait apparaître $A$. L'identité~\eqref{eq:trace} le résume~: la partie
antisymétrique contribue au second ordre, avec un signe opposé à celui de la partie
symétrique. Elle déplace donc le niveau de capital.

\begin{cle}
\textbf{La tension exacte du mémoire.} $A$ ne déplace pas la fluctuation observable
(proposition~\ref{prop:dirichlet}) mais déplace le capital (proposition~\ref{prop:trace}).
Une grandeur qui compte pour la réponse et reste invisible dans la donnée~: c'est la
définition même d'un problème d'\textbf{identification partielle}. On ne peut pas produire
un nombre, on peut produire une bande. L'énumération des $2^{10}=1024$ sommets de la boîte
admissible $|A_{jk}|\le t\,S_{jk}$ n'est donc pas un artifice de calcul, c'est la seule
réponse honnête à la question posée.
\end{cle}

% =====================================================================
\section*{8. La version d'une minute}

\begin{quote}
\small
On peut toujours couper la matrice de contagion en deux~: une partie symétrique, qui dit
quels piliers sont liés, et une partie antisymétrique, qui dit dans quel sens. Si on
renverse le temps, la première ne bouge pas et la seconde change de signe. Or la
fluctuation d'un processus, sa partie martingale, ne dépend que de la première. Donc
n'importe quelle statistique construite sur des co-occurrences annuelles prend la même
valeur sur le processus et sur son renversé, et ne peut pas identifier la direction. Le
placebo à $-0{,}33$ n'est pas un échec de mesure, c'est le théorème. Sauf que la direction,
elle, change le capital, parce qu'elle apparaît dans l'inverse de Leontief. D'où
l'identification partielle~: je ne prétends pas à un nombre, je borne.
\end{quote}

\begin{cle}[colback=orange!8,colframe=orange!70!black]
\textbf{Ce qu'il ne faut pas surinterpréter.} Deux réserves à assumer, dans l'esprit
Vasicheck. \textbf{(i)} $W$ est un noyau de contagion, \emph{pas} littéralement un
générateur de chaîne de Markov. On emprunte l'algèbre de la réversibilité comme cadre de
lecture, on ne pose pas que la cascade \emph{est} une chaîne réversible~: c'est une
lecture, pas une hypothèse supplémentaire. \textbf{(ii)} La forme correcte du courant est
pondérée par $\pi$, et le $S=\tfrac12(W+W\T)$ du mémoire en est le cas particulier
$\pi$ uniforme. L'écart mesuré est faible ($\max_j|\pi_j-1/5|=0{,}069$), donc les deux
lectures coïncident en pratique, mais c'est la version pondérée qui est exacte.
\end{cle}

\smallskip
\sloppy
\noindent\footnotesize\textcolor{navy}{\textbf{Source.}} Vérifications numériques~:
script \texttt{6\_contagion/40\_reversibilite\_martingale.py}, figure
\texttt{Z11\_reversibilite\_martingale.png}. Prolonge la fiche
\emph{réversibilité et martingalisation}, qui donne les valeurs mesurées
($\max|J|=0{,}044$, $\mathcal{E}_P=\mathcal{E}_S$ à $10^{-15}$ près, $\sigma=0{,}070$
contre $5\times10^{-33}$ après symétrisation, $\|A\|/\|W\|=40{,}7\,\%$). Chapitre~9 du
mémoire pour la frontière d'identification, chapitre~10 pour l'énumération des sommets.

\end{document}
